\documentclass[aspectratio=169,11pt]{beamer}

% --- Theme ---
\usetheme{metropolis}
\usepackage{appendixnumberbeamer}

% --- Packages ---
\usepackage{fontspec}
\usepackage{minted}
\setminted{fontsize=\small,breaklines=true}
\usepackage{booktabs}
\usepackage{tikz}
\usepackage{fontawesome5}

% --- Colors ---
\definecolor{healthgreen}{RGB}{34,139,34}
\definecolor{alertred}{RGB}{220,53,69}
\setbeamercolor{alerted text}{fg=alertred}

% --- Metadata ---
\title{Module 3: Data cleaning in health contexts}
\subtitle{Data analysis with Python for health specialists}
\author{Ya\'{e} Ulrich Gaba}
\institute{AIRINA Labs}
\date{2026}

\begin{document}

% ============================================================
\maketitle

% ============================================================
\begin{frame}{Why data cleaning matters in health}
Every health dataset has problems:
\begin{itemize}
  \item Lab values entered with wrong units
  \item Patients appearing twice under different IDs
  \item Dates recorded as \texttt{03/04/2023} --- March 4 or April 3?
  \item Blood glucose of 5000~mg/dL (error) vs.\ 500~mg/dL (real emergency)
\end{itemize}

\vspace{6pt}
Data cleaning is \alert{not optional pre-processing}. It is a \textbf{clinical reasoning task}.

\vspace{6pt}
\begin{alertblock}{Real impact}
In a 2019 EHR study, 18\% of lab values had quality issues. Cleaning decisions changed CKD prevalence estimates by over 3 percentage points.
\end{alertblock}
\end{frame}

% ============================================================
\begin{frame}{Missing data: the three mechanisms}
\begin{table}
\begin{tabular}{lp{4.5cm}p{4.5cm}}
\toprule
\textbf{Type} & \textbf{Example} & \textbf{If you drop} \\
\midrule
\textbf{MCAR} & Lab sample hemolyzed in transport & No bias, less power \\
\textbf{MAR} & HbA1c not ordered for non-diabetics & Bias if you ignore indication \\
\textbf{MNAR} & Sick patients skip follow-up & \alert{Underestimate severity} \\
\bottomrule
\end{tabular}
\end{table}

\vspace{6pt}
The mechanism determines your strategy. \alert{MNAR has no statistical fix} --- only domain knowledge.
\end{frame}

% ============================================================
\begin{frame}[fragile]{Detecting missing values}
\begin{minted}{python}
import pandas as pd
import numpy as np

# Count missing per column
print(df.isnull().sum())

# Percentage missing
print((df.isnull().sum() / len(df) * 100).round(1))

# Rows with any missing value
print(f"Rows with missing data: "
      f"{df.isnull().any(axis=1).sum()} / {len(df)}")
\end{minted}

\vspace{4pt}
\begin{exampleblock}{Visual patterns}
Use \texttt{missingno} for instant visual summaries: \texttt{msno.matrix(df)}, \texttt{msno.bar(df)}, \texttt{msno.heatmap(df)}.
\end{exampleblock}
\end{frame}

% ============================================================
\begin{frame}[fragile]{Imputation strategies}
\begin{columns}
\column{0.48\textwidth}
\textbf{Continuous variables:}
\begin{minted}{python}
# Median (robust to outliers)
df["systolic_bp"] = df["systolic_bp"]\
    .fillna(df["systolic_bp"].median())

# Group-specific median
df["systolic_bp"] = df.groupby(
    "age_group")["systolic_bp"]\
    .transform(
    lambda x: x.fillna(x.median()))
\end{minted}

\column{0.48\textwidth}
\textbf{Categorical variables:}
\begin{minted}{python}
# Mode imputation
df["sex"] = df["sex"].fillna(
    df["sex"].mode()[0])

# KNN imputation (multivariate)
from sklearn.impute import KNNImputer
imputer = KNNImputer(n_neighbors=5)
df[num_cols] = imputer.fit_transform(
    df[num_cols])
\end{minted}
\end{columns}

\vspace{6pt}
\textbf{Rule:} use \alert{median} for skewed lab values, \textbf{mean} for symmetric data, \textbf{mode} for categories.
\end{frame}

% ============================================================
\begin{frame}[fragile]{Handling duplicates}
\begin{minted}{python}
# Exact duplicates (all columns identical)
print(f"Exact duplicates: {patients.duplicated().sum()}")

# Remove exact duplicates
patients_clean = patients.drop_duplicates()

# Keep only first visit per patient
first_visits = patients.drop_duplicates(
    subset="patient_id", keep="first")
\end{minted}

\begin{alertblock}{Clinical context matters}
Same patient, same visit, identical values $\to$ true duplicate. Same patient, different dates $\to$ longitudinal data --- keep all visits!
\end{alertblock}
\end{frame}

% ============================================================
\begin{frame}[fragile]{Standardizing ICD codes and dates}
\begin{minted}{python}
# ICD codes: uppercase, strip, add dot
cleaned = (diagnoses.str.strip().str.upper()
           .str.replace(r"^([A-Z]\d{2})(\d+)$",
                        r"\1.\2", regex=True))

# Dates: parse mixed formats to ISO
parsed = pd.to_datetime(dates, format="mixed",
                        dayfirst=False)
print(parsed.dt.strftime("%Y-%m-%d"))
\end{minted}

\begin{alertblock}{Date ambiguity}
\texttt{03/04/2023} = March 4 (US) or April 3 (Europe). Always check the source. Use \texttt{dayfirst=True} for European dates.
\end{alertblock}
\end{frame}

% ============================================================
\begin{frame}[fragile]{Unit conversions}
\begin{table}
\begin{tabular}{llll}
\toprule
\textbf{Analyte} & \textbf{US} & \textbf{SI} & \textbf{Convert} \\
\midrule
Glucose & mg/dL & mmol/L & $\div$ 18.018 \\
Cholesterol & mg/dL & mmol/L & $\div$ 38.67 \\
Creatinine & mg/dL & \textmu mol/L & $\times$ 88.42 \\
\bottomrule
\end{tabular}
\end{table}

\begin{minted}{python}
# Values > 30 are likely in mg/dL
def standardize_glucose(value):
    if value > 30:
        return value / 18.018  # mg/dL -> mmol/L
    return value  # already mmol/L
\end{minted}
\end{frame}

% ============================================================
\begin{frame}[fragile]{Data validation}
\begin{minted}{python}
# Range checks
ranges = {"age": (0, 120), "systolic_bp": (60, 300),
          "hba1c": (2.0, 20.0), "bmi": (8, 80)}

for col, (low, high) in ranges.items():
    out = df[(df[col] < low) | (df[col] > high)]
    if len(out) > 0:
        print(f"WARNING: {len(out)} values in "
              f"'{col}' outside [{low}, {high}]")

# Cross-field: diastolic < systolic
bp_errors = df[df["diastolic_bp"] >= df["systolic_bp"]]
\end{minted}

\vspace{4pt}
Not every ``impossible'' value is an error. HR of 250~bpm = rare but real (SVT). Glucose of 800~mg/dL = DKA.
\end{frame}

% ============================================================
\begin{frame}[fragile]{Automated quality report}
\begin{minted}{python}
def data_quality_report(df):
    report = []
    for col in df.columns:
        n_miss = df[col].isnull().sum()
        row = {"column": col,
               "dtype": str(df[col].dtype),
               "n_missing": n_miss,
               "pct_missing": round(n_miss/len(df)*100,1),
               "n_unique": df[col].nunique()}
        if df[col].dtype in ["float64", "int64"]:
            row["min"] = df[col].min()
            row["max"] = df[col].max()
        report.append(row)
    return pd.DataFrame(report)
\end{minted}

Run this \alert{before and after} cleaning. Document every decision.
\end{frame}

% ============================================================
\begin{frame}{What you can do after this module}
\begin{enumerate}
  \item Detect missing values and choose the right imputation strategy
  \item Distinguish MCAR, MAR, and MNAR and explain why it matters
  \item Remove true duplicates without losing longitudinal data
  \item Standardize ICD codes, dates, units, and text fields
  \item Validate data against clinically plausible ranges
  \item Generate before/after quality reports to document your pipeline
\end{enumerate}

\vspace{12pt}
\centering
\Large \alert{Next: Module 4 --- Descriptive statistics and epidemiological measures}
\end{frame}

% ============================================================
\begin{frame}{Exercises (Hour 3)}
\begin{enumerate}
  \item \textbf{Quality report:} Generate a data quality report for the provided messy dataset
  \item \textbf{Fix sex and ranges:} Standardize the sex column; set impossible ages and BPs to NaN
  \item \textbf{Unit conversion:} Convert mixed glucose (mg/dL vs.\ mmol/L) and HbA1c (IFCC to NGSP)
  \item \textbf{Mini-project:} Apply the complete cleaning pipeline, compare before/after reports, save cleaned CSV
\end{enumerate}
\end{frame}

\end{document}
